\begin{solution}{56}
    如图 a 所示，光线 1 在上侧柱面 P 点处傍轴垂直入射，入射角为 $\theta$ ，折射角为 $\theta_{0}$ ，由折射定律有
    \begin{equation}
        n \sin \theta = n _ {0} \sin \theta_{0}
        \label{eq:1.s56}
    \end{equation}
    其中 n 和 $n_{0}$ 分别玻璃与空气的折射率。光线在下侧柱面 Q 点处反射，入射角与反射角分别为 i 和 $i'$ ，由反射定律有
    \begin{equation}
        i = i^{\prime}
        \label{eq:2.s56}
    \end{equation}
    光线在下侧柱面 Q 点的反射线交上侧柱面于 P' 点，并由 P' 点向上侧柱面折射，折射光线用 1" 表示；光线 1" 正好与 P' 点处的入射光线 2 的反射光线 2' 相遇，发生干涉。考虑光波反射时的半波损失，光线 1" 与光线 2' 在 P' 点处光程差 $\Delta L$ 为
    \begin{equation}
        \Delta L = \left[ n \left(- z _ {\mathrm{p}}\right) + n _ {0} \left(\mathrm{PQ} + \mathrm{P}^{\prime}\mathrm{Q}\right) + \frac {\lambda}{2} \right] - n \left(- z _ {\mathrm{p}^{\prime}}\right) = n _ {0} \left(\mathrm{PQ} + \mathrm{P}^{\prime}\mathrm{Q}\right) - n \left(z _ {\mathrm{p}} - z _ {\mathrm{p}^{\prime}}\right) + \frac {\lambda}{2}
        \label{eq:3.s56}
    \end{equation}
    式中 $\lambda$ 为入射光线在真空中的波长， $n_{0}=1.00$ 。由题意， $R_{1}$ 和 $R_{2}$ 远大于傍轴光线干涉区域所对应的两柱面间最大间隙；因而在傍轴垂直入射情况下有
    \begin{equation}
        \theta \approx 0, i^{\prime} = i << 1
        \label{eq:4.s56}
    \end{equation}
    \eqref{eq:1.s56}式成为
    \begin{equation}
        n \theta \approx n _ {0} \theta_{0}
        \label{eq:5.s56}
    \end{equation}
    亦即
    \begin{equation}
        1 >> \theta_{0} \approx n \theta \approx 0
        \label{eq:6.s56}
    \end{equation}
    在傍轴条件下，柱面上 P、Q 两处切平面的法线近似平行，因此
    \begin{equation}
        1 >> i^{\prime} = i \approx \theta_{0} \approx 0
        \label{eq:7.s56}
    \end{equation}
    从而，在 P、Q 两处，不仅切平面的法线近似平行，而且在上下表面的反射光线、折射光线均近似平行于入射线，因而也近似平行于 Z 轴，从而 P' 与 P 点近似重合，即
    \begin{equation}
        z _ {\mathrm{p^{\prime}}} \approx z _ {\mathrm{p}}
        \label{eq:8.s56}
    \end{equation}
    且 PQ 近似平行于 Z 轴，因而长度
    \begin{equation}
        \mathrm{P^{\prime}Q} \approx \mathrm{PQ} \approx z _ {\mathrm{p}} - z _ {\mathrm{q}}
        \label{eq:9.s56}
    \end{equation}
    由\eqref{eq:3.s56}\eqref{eq:9.s56}式得
    \begin{equation}
        \Delta L = 2 n _ {0} \cdot \mathrm{PQ} + \frac {\lambda}{2} = 2 n _ {0} \left(z _ {\mathrm{p}} - z _ {\mathrm{Q}}\right) + \frac {\lambda}{2}
        \label{eq:10.s56}
    \end{equation}
    可以将\eqref{eq:10.s56}式右端的 z-坐标近似用 x-或 y-坐标表出。为此，引入一个近似公式。如图 b 所示，设置于平面上的柱面透镜与平面之间的空气隙的厚度为 e，柱面半径为 R。对三边边长分别为 R、$R - e$ 和 $r$ 的直角三角形有
    \begin{equation}
        R^{2} = (R - e)^{2} + r^{2}
        \label{eq:11.s56}
    \end{equation}
    即
    \begin{equation}
        2 R e - e^{2} = r^{2}
        \label{eq:12.s56}
    \end{equation}
    在光线傍轴垂直入射时， $e \ll R$ ，可略去\eqref{eq:12.s56}式左端的 $e^{2}$ ，故
    \begin{equation}
        e = \frac {r^{2}}{2 R}
        \label{eq:13.s56}
    \end{equation}
    在光线傍轴垂直入射时，前面已证近似有 PQ//Z 轴。故可将上、下两个柱面上的 P、Q 两点的坐标取为 $\mathrm{P}(x,y,z_{\mathrm{p}})$ 、 $\mathrm{Q}(x,y,z_{\mathrm{Q}})$ ，如图 c 所示。根据\eqref{eq:13.s56}式可知，P、Q 两点到 XOY 切平面的距离分别为
    \begin{equation}
        e _ {1} = z _ {\mathrm{P}} = \frac {x^{2}}{2 R _ {1}}, e _ {2} = - z _ {\mathrm{Q}} = \frac {y^{2}}{2 R _ {2}}
        \label{eq:14.s56}
    \end{equation}
    最后，光线在上、下两个柱面反射并相遇时，其光程差 $\Delta L$ 为
    \begin{equation}
        \begin{array}{r l}
            \Delta L & \approx 2 n _ {0} \left(z _ {\mathrm{P}} - z _ {\mathrm{Q}}\right) + \frac {\lambda}{2} = 2 n _ {0} \left(e _ {1} + e _ {2}\right) + \frac {\lambda}{2} \\
            & = 2 n _ {0} \left(\frac {x^{2}}{2 R _ {1}} + \frac {y^{2}}{2 R _ {2}}\right) + \frac {\lambda}{2} = n _ {0} \left(\frac {x^{2}}{R _ {1}} + \frac {y^{2}}{R _ {2}}\right) + \frac {\lambda}{2}
        \end{array}
        \label{eq:15.s56}
    \end{equation}
    若 P、Q 两点在 XOY 平面的投影点 $(x, y)$ 落在第 k 级亮(暗)纹上，则 $\Delta L$ 须满足条件
    \begin{equation}
        \Delta L = n _ {0} \left(\frac {x^{2}}{R _ {1}} + \frac {y^{2}}{R _ {2}}\right) + \frac {\lambda}{2} = \left\{
        \begin{array}{l l}
            k \lambda , & k = 1, 2, \dots , \quad \text {亮环} \\
            (k + \frac {1}{2}) \lambda , & k = 0, 1, 2, \dots , \quad \text {暗环}
        \end{array}
        \right.
        \label{eq:16.s56}
    \end{equation}
    \eqref{eq:16.s56}式中亮环条件对应于第 k 级亮纹上的点 $(x, y, z)$ 的 x-、y-坐标满足的方程。

    更具体地，不妨假设 $R_{1} > R_{2}$ ，根据\eqref{eq:16.s56}式中的亮环条件，可得第 k 级亮纹的方程为
    \begin{equation}
        \frac {x^{2}}{A _ {k}^{2}} + \frac {y^{2}}{B _ {k}^{2}} = 1, k = 1, 2 \dots
        \label{eq:17.s56}
    \end{equation}
    它们是椭圆亮环纹，其半长轴与半短轴分别为
    \begin{equation}
        A _ {k} = \sqrt {\left(k - \frac {1}{2}\right) R _ {1} \lambda / n _ {0}}, k = 1, 2, \dots , B _ {k} = \sqrt {\left(k - \frac {1}{2}\right) R _ {2} \lambda / n _ {0}}, k = 1, 2, \dots
        \label{eq:18.s56}
    \end{equation}
\end{solution}